\documentclass[11pt]{article}
% criterion-13: adopted
\usepackage[a4paper,margin=1.1in]{geometry}
\usepackage{amsmath,amssymb,amsthm}
\usepackage[T1]{fontenc}
\usepackage{lmodern}
\emergencystretch=3em

\newtheorem{theorem}{Theorem}
\newtheorem{proposition}{Proposition}
\newtheorem{lemma}{Lemma}
\newtheorem{definition}{Definition}
\newtheorem{remark}{Remark}

\newcommand{\At}{\mathbb{A}_t}
\newcommand{\Ft}{\mathbb{F}_t}
\newcommand{\Nev}{\mathcal N^{\mathrm{ev}}}
\newcommand{\Uev}{U^{\mathrm{hyb,ev}}}
\newcommand{\HC}{\operatorname{HC}}
\newcommand{\Yc}{\check Y}
\newcommand{\Z}{\mathbb Z}
\newcommand{\ev}{\mathrm{ev}}
\newcommand{\qbinom}[2]{\genfrac{[}{]}{0pt}{}{#1}{#2}_{q}}

\title{\textbf{A stand-alone proof of the even hybrid centre theorem for $GL_1$}}
\author{}
\date{}

\begin{document}
\maketitle

\noindent
This note proves the three clauses of \textbf{goalv0a}. It is the
notation-migrated edition of August 10--11, 2026 of the refereed proof of July 19,
2026. The argument is unchanged --- only symbols were renamed, by the Sandbox-wide
dictionary $Q\mapsto q$, $q\mapsto v$, $z\mapsto t$: the relation formerly written
$Q=q^{2}$ now reads $q=v^{2}$, and the family parameter formerly written $z$ is now
written $t$. (In the current Sandbox-wide frame, $v$ and $z$ denote the square roots
$v^{2}=q$ and $z^{2}=t$; only $v$ and $t$ occur below.) Since September 15, 2026 the
rings are printed in blackboard face, $\At$ and $\Ft$ --- a change of face only: the two
macro definitions and nothing else (revised [2026-09-15-12:08 · hypervisorA · Opus 5 Max · castalia:SandboxHQ]). The result identifies the
even hybrid integral form of the family quantum $GL_1$ with its centre and with a
Newton (divided-class) lattice; clause~(iii) is the bilateral quantum P\'olya
theorem of Harman and Hopkins, transported to the multiplicative grid of the
toral algebra (attribution remark at the end).

\section{Conventions}

Throughout, $v$ and $t$ are independent indeterminates and
\[
  \At=\Z[v^{\pm1},t^{\pm1}],\qquad
  \Ft=\mathbb Q(v)[t^{\pm1}],\qquad
  q=v^{2}.
\]
Thus $\At\subseteq\Ft$ is a subring, and $\Z[q^{\pm1}]=\Z[v^{\pm2}]\subseteq
\Z[v^{\pm1}]\subseteq\At$.

\emph{The algebra $U$.}\;
$U$ is the $\Ft$-algebra with invertible pairwise-commuting generators
$K_{n,+},K_{n,-}$ ($n\in\Z$) subject to
$K_{0,\pm}=1$, $K_{m+n,\pm}=K_{m,\pm}K_{n,\pm}$ and $K_{n,+}K_{n,-}=t^{\,n}$.
Writing $T^{\pm}=K_{1,\pm}$, the monomials $(T^{+})^{w}$ ($w\in\Z$) form a
$\Ft$-basis; $U$ is commutative. A \emph{toral} element is an element of
$\Ft[(T^{+})^{\pm1}]$; it is \emph{even} when it is a Laurent polynomial in
\[
  \Yc:=(T^{+})^{2}.
\]
Even toral elements thus form the Laurent-polynomial algebra
$\Ft[\Yc^{\pm1}]$.

\emph{Evaluations.}\;
For $\chi\in\Z$, let $\ev_\chi\colon\Ft[\Yc^{\pm1}]\to\Ft$ be the
$\Ft$-algebra homomorphism substituting the unit $q^{\chi}\in\At$ for $\Yc$,
i.e.
\[
  \ev_\chi\Bigl(\textstyle\sum_{u}c_u\Yc^{u}\Bigr)=\sum_u c_u\,q^{u\chi}
  \qquad(c_u\in\Ft).
\]
It is the unique $\At$-linear ring homomorphism with $\ev_\chi(\Yc)=q^{\chi}$,
and it fixes $\Ft$ (hence $\At$) pointwise. The \emph{even Newton lattice} is
\[
  \Nev=\bigl\{\,f\in\Ft[\Yc^{\pm1}]:\ \ev_\chi(f)\in\At\ \text{for every }
  \chi\in\Z\,\bigr\}.
\]
Because $\Phi^{+}=\emptyset$ at $GL_1$, the \emph{even hybrid family integral
form} is $\Uev=\Nev\subseteq U$, with central part
$Z(\Uev)=\Uev\cap Z\bigl(U\otimes_{\Ft}\operatorname{Frac}(\At)\bigr)$. Every
element of $U$ is toral, so the Harish--Chandra projection $\HC$ is the identity
map on $U$; the Weyl group is $W=\{1\}$, acting trivially, so
$(\Nev)^{W}=\Nev$.

\emph{Newton divided classes.}\;
For $r\ge0$,
\[
  \nu_{r}(X)=\prod_{s=0}^{r-1}\frac{X-q^{s}}{q^{r}-q^{s}}\ \in\ \mathbb Q(v)[X],
  \qquad \nu_0(X)=1 .
\]
Each $\nu_r$ has degree exactly $r$ (its leading coefficient
$\prod_{s=0}^{r-1}(q^{r}-q^{s})^{-1}$ is a nonzero element of $\mathbb Q(v)$),
so $\{\nu_0,\dots,\nu_d\}$ is a $\Ft$-basis of the module
$\Ft[\Yc]_{\le d}$ of even toral polynomials of degree $\le d$.

\medskip
We write $\Z[q^{\pm1}]$-\emph{integral} for ``lies in $\Z[q^{\pm1}]$''.
For $m\in\Z$ and $r\ge0$ put
\[
  B_r(m):=\nu_r(q^{m})=\prod_{s=0}^{r-1}\frac{q^{m}-q^{s}}{q^{r}-q^{s}},
  \qquad B_0(m)=1 .
\]

\section{Clause (i): $\Uev$ is an $\At$-subalgebra}

\begin{proposition}
$\Uev=\Nev$ is an $\At$-subalgebra of $U$. (This is the bilateral,
$t$-spectator analogue of the fact that quantum integer-valued polynomials
form a ring \cite[\S1]{HarmanHopkins}.)
\end{proposition}

\begin{proof}
For each $\chi\in\Z$ the map $\ev_\chi\colon\Ft[\Yc^{\pm1}]\to\Ft$ is a ring
homomorphism, and $\At\subseteq\Ft$ is a subring; hence the preimage
$\ev_\chi^{-1}(\At)$ is a subring of $\Ft[\Yc^{\pm1}]$. It contains $\At$,
because $\ev_\chi$ fixes $\At$ pointwise, so $\ev_\chi^{-1}(\At)$ is an
$\At$-subalgebra. By definition
\[
  \Nev=\bigcap_{\chi\in\Z}\ev_\chi^{-1}(\At)
\]
is an intersection of $\At$-subalgebras of $\Ft[\Yc^{\pm1}]$, hence itself an
$\At$-subalgebra (it contains $1$ since $\ev_\chi(1)=1\in\At$, and is closed
under sums, products and $\At$-scalars). As $\Ft[\Yc^{\pm1}]\subseteq U$, this
is an $\At$-subalgebra of $U$.
\end{proof}

\section{Clause (ii): the Harish--Chandra isomorphism is the identity}

\begin{proposition}
The Harish--Chandra projection restricts to an isomorphism of $\At$-algebras
$\HC\colon Z(\Uev)\xrightarrow{\ \sim\ }(\Nev)^{W}$, and this isomorphism is the
identity map of $\Nev$. (The rank-one, $\Phi^{+}=\emptyset$ degeneration of the
even hybrid Harish--Chandra isomorphism established at $GL_2$ in
\cite{SandboxGL2}; the classical ancestor is the Harish--Chandra isomorphism
itself, whose quantum-group form the $GL_2$ record cites.)
\end{proposition}

\begin{proof}
The algebra $U$ is commutative, being spanned by monomials in the pairwise
commuting invertible generators $K_{n,\pm}$. Base change is compatible with
commutativity, so $U\otimes_{\Ft}\operatorname{Frac}(\At)$ is a commutative ring
and equals its own centre. Therefore
\[
  Z(\Uev)=\Uev\cap Z\bigl(U\otimes_{\Ft}\operatorname{Frac}(\At)\bigr)
        =\Uev\cap\bigl(U\otimes_{\Ft}\operatorname{Frac}(\At)\bigr)
        =\Uev=\Nev,
\]
the middle equality because $U\hookrightarrow U\otimes_{\Ft}\operatorname{Frac}(\At)$
is injective --- $\Ft\hookrightarrow\operatorname{Frac}(\At)=\mathbb Q(v,t)$ and $U$ is
free over $\Ft$ --- so $\Uev=\Nev\subseteq U$ is literally a subset of the base-changed
algebra.
Every element of $U$ is toral, so $\HC=\mathrm{id}_U$, and its restriction to
$Z(\Uev)=\Nev$ is $\mathrm{id}_{\Nev}$. Finally $W=\{1\}$ acts trivially, so
$(\Nev)^{W}=\Nev$. Hence
\[
  \HC\colon Z(\Uev)=\Nev\ \xrightarrow{\ \mathrm{id}\ }\ \Nev=(\Nev)^{W}
\]
is the identity map of the $\At$-algebra $\Nev$ (an $\At$-algebra by clause~(i)),
which is manifestly an isomorphism of $\At$-algebras.
\end{proof}

\noindent
There is no content here beyond commutativity: at $GL_1$ the rational centre is
everything, $\HC$ is the identity, and Weyl invariance is vacuous.

\section{Clause (iii): the Newton presentation}

Write
\[
  S\ :=\ \sum_{u\in\Z}\ \sum_{r\ge0}\ \At\,\Yc^{\,u}\,\nu_{r}(\Yc)\ \subseteq\
  \Ft[\Yc^{\pm1}]
\]
for the $\At$-submodule generated by the shifted divided classes
$\Yc^{u}\nu_r(\Yc)$.

\begin{theorem}[Even Newton presentation for $GL_1$]\label{thm:iii}
$\displaystyle \Nev=S=\sum_{u\in\Z}\sum_{r\ge0}\At\,\Yc^{\,u}\,\nu_r(\Yc).$
(The bilateral quantum P\'olya theorem of Harman--Hopkins
\cite[\S\S1,4]{HarmanHopkins}, transported to the multiplicative grid; the
precise dictionary is the Attribution remark below.)
\end{theorem}

The proof occupies the rest of this section: the grid lemma, the forward
inclusion $S\subseteq\Nev$, the windowed reverse inclusion for polynomials, and
the Laurent-clearing step.

\subsection{The grid values are $\Z[q^{\pm1}]$-integral}

\begin{lemma}\label{lem:grid}
For all $m\in\Z$ and $r\ge0$ one has $B_r(m)=\nu_r(q^{m})\in\Z[q^{\pm1}]$
(the Gaussian-binomial value pattern of \cite[\S1]{HarmanHopkins}, extended
to negative nodes as their \S4 bilateral localization requires).
Explicitly:
\begin{enumerate}
\item[(a)] if $0\le m<r$ then $B_r(m)=0$;
\item[(b)] if $m\ge r\ge0$ then $B_r(m)=\qbinom{m}{r}\in\Z[q]$;
\item[(c)] if $m=-n$ with $n\ge1$ then
  \[
    B_r(-n)=(-1)^{r}\,q^{-nr-\binom{r}{2}}\,\qbinom{n+r-1}{r}\ \in\ \Z[q^{\pm1}].
  \]
\end{enumerate}
\end{lemma}

\begin{proof}
\emph{(a)}\; If $0\le m<r$, the factor of the product at $s=m$ has numerator
$q^{m}-q^{m}=0$, so $B_r(m)=0$.

\emph{(b)}\; Assume $m\ge r\ge0$. Pulling $q^{s}$ out of each numerator and
denominator factor,
\[
  B_r(m)=\prod_{s=0}^{r-1}\frac{q^{s}(q^{m-s}-1)}{q^{s}(q^{r-s}-1)}
        =\prod_{s=0}^{r-1}\frac{q^{m-s}-1}{q^{r-s}-1}.
\]
Re-indexing by $i=r-s$ (so $i$ runs from $1$ to $r$),
\[
  B_r(m)=\prod_{i=1}^{r}\frac{q^{(m-r)+i}-1}{q^{i}-1}=\qbinom{m}{r},
\]
the Gaussian binomial coefficient. It lies in $\Z[q]$: this is standard from the
$q$-Pascal recurrence $\qbinom{m}{r}=\qbinom{m-1}{r-1}+q^{r}\qbinom{m-1}{r}$ with
$\qbinom{m}{0}=1$, by induction on $m$.

\emph{(c)}\; Let $m=-n$, $n\ge1$. In each factor write the numerator as
$q^{-n}-q^{s}=-q^{-n}(q^{n+s}-1)$ and the denominator as
$q^{r}-q^{s}=q^{s}(q^{r-s}-1)$. Then
\[
  B_r(-n)=\prod_{s=0}^{r-1}\frac{-q^{-n}(q^{n+s}-1)}{q^{s}(q^{r-s}-1)}
        =(-1)^{r}\,q^{-nr}\,q^{-\sum_{s=0}^{r-1}s}
          \prod_{s=0}^{r-1}\frac{q^{n+s}-1}{q^{r-s}-1}.
\]
Here $\sum_{s=0}^{r-1}s=\binom{r}{2}$. For the remaining product, the
denominators give $\prod_{s=0}^{r-1}(q^{r-s}-1)=\prod_{i=1}^{r}(q^{i}-1)$ and the
numerators give $\prod_{s=0}^{r-1}(q^{n+s}-1)=\prod_{j=n}^{n+r-1}(q^{j}-1)$, so
\[
  \prod_{s=0}^{r-1}\frac{q^{n+s}-1}{q^{r-s}-1}
  =\frac{\prod_{j=n}^{n+r-1}(q^{j}-1)}{\prod_{i=1}^{r}(q^{i}-1)}
  =\frac{[n+r-1]_q!\,/\,[n-1]_q!}{[r]_q!}=\qbinom{n+r-1}{r},
\]
where $[k]_q=(q^{k}-1)/(q-1)$ and the common powers of $(q-1)$ cancel. The
Gaussian binomial $\qbinom{n+r-1}{r}\in\Z[q]$ by part (b). Collecting the
prefactor gives
\[
  B_r(-n)=(-1)^{r}\,q^{-nr-\binom{r}{2}}\,\qbinom{n+r-1}{r}\in\Z[q^{\pm1}].
  \qedhere
\]
\end{proof}

\subsection{Forward inclusion $S\subseteq\Nev$}

\begin{lemma}\label{lem:fwd}
For every $u\in\Z$ and $r\ge0$, $\Yc^{u}\nu_r(\Yc)\in\Nev$. Hence $S\subseteq\Nev$.
(The transported forward half of \cite[\S1]{HarmanHopkins}: $q$-binomials are
quantum integer-valued.)
\end{lemma}

\begin{proof}
Fix $u,r$ and $\chi\in\Z$. Since $\ev_\chi$ is a ring homomorphism with
$\ev_\chi(\Yc)=q^{\chi}$,
\[
  \ev_\chi\bigl(\Yc^{u}\nu_r(\Yc)\bigr)=q^{u\chi}\,\nu_r(q^{\chi})
  =q^{u\chi}\,B_r(\chi).
\]
By Lemma~\ref{lem:grid}, $B_r(\chi)\in\Z[q^{\pm1}]\subseteq\At$, and $q^{u\chi}$
is a unit of $\At$; so the product lies in $\At$. As $\chi$ was arbitrary,
$\Yc^{u}\nu_r(\Yc)\in\Nev$. The lattice $\Nev$ is an $\At$-module (indeed an
$\At$-algebra, clause~(i)), so it contains every $\At$-combination of these
generators, i.e.\ $S\subseteq\Nev$.
\end{proof}

\subsection{Reverse inclusion for polynomials: windowed Newton interpolation}

\begin{lemma}\label{lem:rev}
Let $g\in\Nev$ be an even toral \emph{polynomial}, $g\in\Ft[\Yc]$, of degree
$\le d$. Then $g=\sum_{r=0}^{d}c_r\,\nu_r(\Yc)$ with all $c_r\in\At$. In
particular $g\in\sum_{r\ge0}\At\,\nu_r(\Yc)\subseteq S$.
(P\'olya's windowed-interpolation argument as $q$-deformed in the proof of
\cite[\S1]{HarmanHopkins}, run on the multiplicative grid with $t$ a
spectator.)
\end{lemma}

\begin{proof}
Because $\{\nu_0,\dots,\nu_d\}$ is a $\Ft$-basis of $\Ft[\Yc]_{\le d}$, there are
unique $c_0,\dots,c_d\in\Ft$ with
\[
  g=\sum_{r=0}^{d}c_r\,\nu_r(\Yc).
\]
Evaluate on the window $\chi=m$ for $m=0,1,\dots,d$. As $\ev_m$ is
$\Ft$-linear and $\ev_m(\nu_r(\Yc))=B_r(m)$,
\[
  \ev_m(g)=\sum_{r=0}^{d}B_r(m)\,c_r
  =\sum_{r=0}^{d}V[m,r]\,c_r,
  \qquad V[m,r]:=B_r(m)\ \ (0\le m,r\le d).
\]
By Lemma~\ref{lem:grid}, every entry $V[m,r]\in\Z[q^{\pm1}]$; moreover $V$ is
\emph{lower unitriangular}: $V[m,r]=B_r(m)=0$ for $m<r$ (part (a)), and
$V[m,m]=B_m(m)=\prod_{s=0}^{m-1}\frac{q^{m}-q^{s}}{q^{m}-q^{s}}=1$. Hence
$\det V=1$, and the inverse $V^{-1}$ is again lower unitriangular with entries in
$\Z[q^{\pm1}]$ (back-substitution never leaves $\Z[q^{\pm1}]$, and by uniqueness
this inverse coincides with the inverse taken over the larger ring $\Ft$).

Writing $\mathbf c=(c_0,\dots,c_d)^{\mathsf T}$ and
$\mathbf v=(\ev_0(g),\dots,\ev_d(g))^{\mathsf T}$, we have $\mathbf v=V\mathbf c$
in $\Ft^{d+1}$, hence
\[
  \mathbf c=V^{-1}\mathbf v .
\]
Since $g\in\Nev$, each $\ev_m(g)\in\At$, so $\mathbf v\in\At^{d+1}$; and
$V^{-1}$ has entries in $\Z[q^{\pm1}]\subseteq\At$. Therefore
$\mathbf c=V^{-1}\mathbf v\in\At^{d+1}$, i.e.\ every $c_r\in\At$.

\smallskip
\emph{The $t$-spectator argument.}\;
The map $\ev_m$ substitutes only $\Yc\mapsto q^{m}=v^{2m}$; it does not touch $v$
or $t$. Thus $t$ never enters the matrix $V$, whose entries lie in
$\Z[q^{\pm1}]$. Concretely, expand each coordinate along the free
$\Z[v^{\pm1}]$-basis of $t$-powers: $c_r=\sum_{k}c_{r,k}\,t^{k}$ with
$c_{r,k}\in\mathbb Q(v)$, so that
$\ev_m(g)=\sum_{k}\bigl(\sum_{r}B_r(m)\,c_{r,k}\bigr)t^{k}$. Membership
$\ev_m(g)\in\At=\Z[v^{\pm1},t^{\pm1}]$ means that for each fixed $k$ the
$t^{k}$-coefficient $\sum_r B_r(m)c_{r,k}$ lies in $\Z[v^{\pm1}]$ for all
$m\in\{0,\dots,d\}$. For each fixed $k$ this is precisely the unitriangular
system above, over $\Z[q^{\pm1}]$ with right-hand sides in $\Z[v^{\pm1}]$; it
gives $c_{r,k}\in\Z[v^{\pm1}]$. Hence
$c_r=\sum_k c_{r,k}t^{k}\in\Z[v^{\pm1},t^{\pm1}]=\At$, as claimed. This is the
same conclusion reached above; we spell out the coefficient-wise reading to make
the spectator role of $t$ explicit.

\smallskip
\emph{Base change $\Z[q^{\pm1}]\to\Z[v^{\pm1}]\to\At$.}\;
Two rings appear: the matrix $V$ and its inverse live over
$\Z[q^{\pm1}]=\Z[v^{\pm2}]$, whereas the coordinates $c_r$ and the values
$\ev_m(g)$ live over $\Ft=\mathbb Q(v)[t^{\pm1}]$ --- honest $v$ (odd powers) may
occur, so the $c_r$ need not lie in $\mathbb Q(q)$. This causes no loss: the
extension $\Z[q^{\pm1}]\hookrightarrow\Z[v^{\pm1}]$ is free with basis
$\{1,v\}$, hence flat, and multiplication by a $\Z[q^{\pm1}]$-entry preserves
each of the two grades (even/odd powers of $v$). The reverse computation
$\mathbf c=V^{-1}\mathbf v$ therefore runs identically in each grade, so
integrality of the values over $\Z[v^{\pm1},t^{\pm1}]$ forces integrality of the
coordinates over the same ring. Equivalently, and most directly:
$V^{-1}$ has entries in $\Z[q^{\pm1}]\subseteq\At$ and $\mathbf v\in\At^{d+1}$,
so $\mathbf c=V^{-1}\mathbf v\in\At^{d+1}$.
\end{proof}

\subsection{Laurent clearing and conclusion}

\begin{proof}[Proof of Theorem~\ref{thm:iii}]
The inclusion $S\subseteq\Nev$ is Lemma~\ref{lem:fwd}. For $\Nev\subseteq S$,
let $f\in\Nev\subseteq\Ft[\Yc^{\pm1}]$ be arbitrary. Choose $N\ge0$ so large that
$g:=\Yc^{N}f\in\Ft[\Yc]$ is an honest polynomial (take $N$ at least the largest
magnitude of a negative $\Yc$-exponent occurring in $f$). For every $\chi\in\Z$,
using that $\ev_\chi$ is a ring homomorphism,
\[
  \ev_\chi(g)=\ev_\chi(\Yc^{N})\,\ev_\chi(f)=q^{N\chi}\,\ev_\chi(f).
\]
Here $q^{N\chi}$ is a unit of $\At$ and $\ev_\chi(f)\in\At$, so
$\ev_\chi(g)\in\At$; thus $g\in\Nev\cap\Ft[\Yc]$. By Lemma~\ref{lem:rev},
$g=\sum_{r}c_r\,\nu_r(\Yc)$ with $c_r\in\At$. Multiplying back by the unit
$\Yc^{-N}$,
\[
  f=\Yc^{-N}g=\sum_{r\ge0}c_r\,\Yc^{-N}\nu_r(\Yc)
   \ \in\ \sum_{u\in\Z}\sum_{r\ge0}\At\,\Yc^{\,u}\,\nu_r(\Yc)=S,
\]
each summand being an $\At$-multiple of the generator $\Yc^{u}\nu_r(\Yc)$ with
$u=-N$. Hence $\Nev\subseteq S$, and with Lemma~\ref{lem:fwd}, $\Nev=S$.
\end{proof}

\section{Attribution}

\begin{remark}
Clause~(iii) is the quantum P\'olya theorem of Harman and Hopkins,
\emph{Quantum integer-valued polynomials} (arXiv:1601.06110), in bilateral
Laurent form. Their \S1 Proposition (Prop.~\S1.2, \texttt{prop:qbinom}), the
$q$-analogue of P\'olya's 1919 theorem, states that the ring
$\mathcal R_q^{+}=\{P\in\mathbb Q(q)[x]:P([n]_q)\in\Z[q]\ \forall n\in\mathbb N\}$
is freely generated as a $\Z[q]$-module by the $q$-binomial coefficient
polynomials $\begin{bmatrix}x\\k\end{bmatrix}_{q}$ (the $q$-deformed Newton
basis on the \emph{additive} grid $x=[n]_q$, with denominator
$q^{\binom{k}{2}}[k]_q!$ --- not the classical P\'olya $\binom{x}{k}$); their \S4
Proposition (Prop.~\S4.3, \texttt{prop:localization}),
$\mathcal R_q=\mathcal R_q^{+}\otimes_{\Z[q]}\Z[q^{\pm1}]$, is the bilateral
localization. In the current notation their parameter $q$ coincides with our
$q$ ($=v^{2}$), so their displayed formulas read verbatim on our side.

The node-preserving affine substitution $\Yc=1+(q-1)x$ is exactly their \S4
$z$-variable; it carries the grid node $x=[n]_q$ to $\Yc=q^{n}$ and carries each
basis term $\begin{bmatrix}x\\k\end{bmatrix}_{q}$ onto $\nu_k(\Yc)$ \emph{term by
term} --- the factors $(q-1)$ cancel in every Newton ratio. (They remark further
that on adjoining square roots $K^{2}=z$, $v^{2}=q$ --- which is precisely the
passage to our $T^{+}$ and our $v$ --- the ring $\mathcal R_q[K,v]$ is equivalent
to the Cartan part of Lusztig's integral form of $U_v(\mathfrak{sl}_2)$. Nothing
below rests on that remark.) The values-in-$\Z[q^{\pm1}]$
side matches values-in-$\At$ with $t$ a spectator, since evaluation is
$t$-coefficient-wise and $\mathbb Q(q)\cap\Z[v^{\pm1}]=\Z[q^{\pm1}]$.

The two-variable ($GL_2$) form of this equivalence is established in
\cite{SandboxGL2}: there the grid criterion agrees with membership in the
$\At$-span of the product Newton generators, and passes to the full Laurent
even lattice. The $GL_1$ statement above is the rank-one specialization of
that transport.
\end{remark}

\begin{thebibliography}{9}

\bibitem{HarmanHopkins}
N.~Harman and S.~Hopkins,
\emph{Quantum integer-valued polynomials},
arXiv:1601.06110 (2016).

\bibitem{SandboxGL2}
SandboxA,
\emph{The centre of the even hybrid family quantum $GL_2$}
(closure record of the signed goal, revision A1+R5),
Sandbox record, 2026.

\end{thebibliography}

\end{document}
